\documentclass[twoside, 10pt, a4paper]{article}
\usepackage[utf8]{inputenc}
\usepackage[T1]{fontenc}
\usepackage[portuguese]{babel}

% --- FONTE PADRÃO (Estilo Didático Encorpado) ---
\usepackage{helvet} 
\renewcommand{\familydefault}{\sfdefault}

% --- GEOMETRIA E ESPAÇAMENTO ---
\usepackage[top=1.6cm, bottom=1.5cm, left=0.9cm, right=0.9cm, headheight=22pt, headsep=0.4cm]{geometry}
\usepackage{microtype} 
\usepackage{setspace}
\setstretch{1.0}
\usepackage{parskip}
\setlength{\parskip}{2pt}
\setlength{\parindent}{0pt}

\newlength{\espacoPosTitulo}\setlength{\espacoPosTitulo}{0.3cm}
\newlength{\espacoPreQuestao}\setlength{\espacoPreQuestao}{0.1cm}
\newlength{\espacoQuestao}\setlength{\espacoQuestao}{0.2cm}
\newlength{\espacoAlt}\setlength{\espacoAlt}{0.2cm}

\usepackage{enumitem}
\setlist{itemsep=0.4cm, topsep=0.1cm, parsep=0pt, partopsep=0pt}
\newcommand{\larguraTikz}{0.85\linewidth} 

% --- MATEMÁTICA E CORES ---
\usepackage{amsmath, amsfonts, amssymb, nccmath, mathastext}
\usepackage{xcolor}
\definecolor{indigo}{HTML}{4F46E5}
\definecolor{CorTitUm}{HTML}{FF6F61}
\definecolor{CorTitDois}{HTML}{FF6F61}
\definecolor{CorTitTres}{HTML}{658894}
\definecolor{CorLinha}{HTML}{B0B0B0} 
\definecolor{metal}{HTML}{7F8C8D}
\definecolor{vidro}{HTML}{3498DB}
\definecolor{ouro}{HTML}{F1C40F}
\definecolor{concreto}{HTML}{95A5A6}
\definecolor{madeira}{HTML}{D35400}
\definecolor{CinzaEscuro}{HTML}{4A4A4A} 
\definecolor{CinzaMedio}{HTML}{848484} 
\definecolor{Cinzablack}{HTML}{353535} 

% --- MINI-CABEÇALHO E RODAPÉ AUTORAL (TWOSIDE) ---
\usepackage{fancyhdr}
\pagestyle{fancy}
\fancyhf{} 
\renewcommand{\headrulewidth}{0.3pt} 
\renewcommand{\footrulewidth}{0.3pt}
\fancyhead[LE]{\small\sffamily\bfseries\color{gray!75} DECIMAIS, POTENCIAÇÃO E RADICIAÇÃO}
\fancyhead[RO]{\small\sffamily\color{gray!75} \texttt{profraphaelpaulino.com.br}}
\fancyfoot[LE]{\small \textbf{\thepage}}
\fancyfoot[RO]{\small \textbf{\thepage}}

% --- CAIXAS DE DESTAQUE ---
\usepackage[most]{tcolorbox}
\newtcolorbox{boxresponda}{colback=white, colframe=gray!45, boxrule=0pt, leftrule=1.7pt, sharp corners, enlarge left by=0.4cm, width=\linewidth-0.4cm, left=8pt, right=0pt, top=2pt, bottom=2pt, fontupper=\small}
\definecolor{CorDicaBorda}{HTML}{17c1be} 
\definecolor{CorDicaFundo}{HTML}{f3fbfb} 
\newtcolorbox{boxdica}{colback=CorDicaFundo, colframe=CorDicaBorda, boxrule=0pt, leftrule=2.5pt, sharp corners, width=\linewidth, left=8pt, right=8pt, top=6pt, bottom=6pt, halign=center, fontupper=\small}

% --- TÍTULOS ---
\usepackage{titlesec}
\titleformat{\section}{\color{black}\large\bfseries}{\thesection}{0em}{\vspace{0.1cm}\titlerule[0.8pt]}
\titlespacing*{\section}{0pt}{10pt}{6pt} 
\titleformat{\subsubsection}[runin]{\color{black}\bfseries}{}{0em}{}
\titlespacing*{\subsubsection}{0pt}{0pt}{0.15cm}

% --- COLUNAS E GRÁFICOS ---
\usepackage{multicol}
\setlength{\columnsep}{1cm}
\setlength{\columnseprule}{0.5pt}
\def\columnseprulecolor{\color{CorLinha}}
\usepackage{adjustbox, tikz, pgfplots, circuitikz}
\usetikzlibrary{shadows, shapes.misc, positioning, calc}
\pgfplotsset{compat=1.18}

\begin{document}
\thispagestyle{empty}
\color{Cinzablack}
\vspace*{-1.8cm}

% --- CABEÇALHO PRINCIPAL AUTORAL (Apenas 1ª página) ---
\noindent
\begin{minipage}{\textwidth}
    \fontfamily{phv}\selectfont
    \noindent
    \begin{minipage}[t]{0.70\linewidth}
        {\Large \bfseries \MakeUppercase{Caderno de Atividades}}\\[0.1cm]
        {\large \textmd{\textcolor{CinzaEscuro}{Unidade: Decimais, Potenciação e Radiciação}}}
    \end{minipage}%
    \begin{minipage}[t]{0.30\linewidth}
        \raggedleft
        \small \textbf{Prof. Raphael Paulino}\\
        \textsf{\small \textcolor{indigo}{\texttt{profraphaelpaulino.com.br}}}
    \end{minipage}
    
    \vspace{0.15cm}
    \noindent\rule{\linewidth}{1.2pt}
    \vspace{-0.1cm}
    \footnotesize\textsf{\textbf{ÁREA:} MATEMÁTICA \hfill \textbf{NÍVEL:} 7º ANO}
    \vspace{0.1cm}
    \noindent\rule{\linewidth}{0.4pt}
\end{minipage}

\par\vspace{0.3cm} 
\raggedcolumns
\begin{multicols*}{2}
\vspace{0.1cm}

\subsection*{Capítulo 1: Representação Decimal}
\vspace{-0.2cm}\noindent\rule{\linewidth}{1.5pt} 
\par\vspace{\espacoPosTitulo}
\textbf{1.} A leitura de instrumentos de precisão, como escalas em maquinários industriais, exige a correta interpretação das subdivisões decimais. O esquema ilustra a ampliação do mostrador de uma máquina de corte a laser, onde o marcador \textbf{P} indica a espessura exata da chapa de aço a ser utilizada.

\begin{center}
\begin{adjustbox}{max width=\linewidth}
\begin{tikzpicture}
% --- APLICANDO DROP SHADOW E CORES HEX ---
\definecolor{azulprincipal}{HTML}{2980B9}
\definecolor{fundocaixa}{HTML}{ECF0F1}
\definecolor{sombracaixa}{HTML}{BDC3C7}
\definecolor{destaque}{HTML}{E74C3C}

% LAYER 1: Sombra e Fundo expandidos (Margens Seguras)
\fill[sombracaixa, rounded corners=4pt, drop shadow={opacity=0.2, shadow xshift=3pt, shadow yshift=-3pt}] (-1, -0.5) rectangle (11, 2.5);
\fill[fundocaixa, rounded corners=4pt] (-1, -0.5) rectangle (11, 2.5);

% LAYER 2: Eixo Numérico
\draw[thick, -latex, azulprincipal] (0.5, 1) -- (9.5, 1);

% Marcas principais e secundárias
\foreach \x in {1, 1.5, 2, 2.5, 3, 3.5, 4, 4.5, 5, 5.5, 6, 6.5, 7, 7.5, 8, 8.5, 9} {
    \draw[thick, azulprincipal] (\x, 0.8) -- (\x, 1.2);
}

% Rótulos com fill=fundocaixa e inner sep para não cruzar linhas
\node[below, fill=fundocaixa, inner sep=2pt, font=\bfseries] at (1, 0.7) {1,20};
\node[below, fill=fundocaixa, inner sep=2pt, font=\bfseries] at (5, 0.7) {1,24};
\node[below, fill=fundocaixa, inner sep=2pt, font=\bfseries] at (9, 0.7) {1,28};

% Marcador P
\fill[destaque] (7.5, 1) circle (0.15);
\node[above, text=destaque, fill=fundocaixa, inner sep=2pt, font=\bfseries] at (7.5, 1.4) {P};
\end{tikzpicture}
\end{adjustbox}
\end{center}

Determine o valor decimal exato correspondente ao ponto \textbf{P} e, em seguida, converta esse valor para a forma de fração irredutível. Justifique seu raciocínio demonstrando a escala adotada no eixo.
\par\vspace{\espacoQuestao}

\noindent\textcolor{CorLinha}{\rule{\linewidth}{0.5pt}} 
\par\vspace{\espacoPreQuestao}
\textbf{2.} A transformação entre frações e números decimais obedece a regras de equivalência de base. Analise as afirmativas a seguir sobre a equivalência e a leitura geométrica dessas representações.

Escreva V para verdadeiro ou F para falso em cada item e justifique matematicamente as afirmativas que julgar falsas.
\begin{enumerate}[label=\Roman*.]
    \item O número decimal correspondente à fração $ \mfrac{15}{4} $ é exatamente \textbf{3,75}, que representa a soma de três inteiros com três quartos.
    \item A leitura correta do número \textbf{0,045} é "quarenta e cinco centésimos".
    \item Transformando o número \textbf{2,8} em fração irredutível, obtemos $ \mfrac{14}{5} $.
\end{enumerate}
\par\vspace{\espacoQuestao}

\noindent\textcolor{CorLinha}{\rule{\linewidth}{0.5pt}} 
\par\vspace{\espacoPreQuestao}
\textbf{3.} Na final olímpica dos 100 metros rasos masculinos, a precisão eletrônica é fundamental para definir as medalhas. O placar oficial registrou os seguintes tempos (em segundos) para os quatro primeiros colocados: Atleta A (\textbf{9,805}), Atleta B (\textbf{9,85}), Atleta C (\textbf{9,089}) e Atleta D (\textbf{9,8}).

Organize os atletas do primeiro ao quarto lugar e calcule a diferença de tempo, em milésimos de segundo, entre o vencedor da prova e o medalhista de bronze.
\par\vspace{\espacoQuestao}

\noindent\textcolor{CorLinha}{\rule{\linewidth}{0.5pt}} 
\par\vspace{\espacoPreQuestao}
\textbf{4.} No projeto de um motor de combustão, as peças importadas geralmente possuem medidas em polegadas fracionárias, enquanto o sistema nacional exige a conversão para o sistema decimal. Um engenheiro precisa ajustar uma porca com especificação de $ \mfrac{7}{16} $ de polegada e uma arruela de $ \mfrac{3}{8} $ de polegada.

Determine o valor decimal exato correspondente a cada uma das peças. Qual delas possui o maior diâmetro interno? Demonstre o cálculo da divisão.
\par\vspace{\espacoQuestao}

\noindent\textcolor{CorLinha}{\rule{\linewidth}{0.5pt}} 
\par\vspace{\espacoPreQuestao}
\textbf{5.} Laboratórios químicos utilizam beckers de precisão para preparar soluções. A graduação lateral do recipiente cilíndrico ilustra o volume ocupado por um reagente azul. A marcação superior máxima corresponde exatamente a \textbf{1 L} (um litro).

\begin{center}
\begin{adjustbox}{max width=\linewidth}
\begin{tikzpicture}
% --- APLICANDO DROP SHADOW E CORES HEX ---
\definecolor{vidro}{HTML}{7F8C8D}
\definecolor{liquido}{HTML}{3498DB}
\definecolor{marcadores}{HTML}{2C3E50}

% LAYER 1: Fundo expandido e Sombra (Margens Seguras para não cortar os textos)
\fill[black!10, rounded corners=4pt, drop shadow={opacity=0.2, shadow xshift=3pt, shadow yshift=-3pt}] (-0.5, -0.5) rectangle (6.5, 7.0);
\fill[white, rounded corners=4pt] (-0.5, -0.5) rectangle (6.5, 7.0);

% LAYER 2: Líquido
\fill[liquido, opacity=0.85, rounded corners=2pt] (1.1, 0.1) rectangle (3.9, 4.3);

% LAYER 3: Vidro e Borda
\draw[ultra thick, vidro, rounded corners=4pt] (1, 6) -- (1, 0) -- (4, 0) -- (4, 6);
\draw[thick, vidro] (1, 6) ellipse (1.5 and 0.3);

% LAYER 4: Marcadores de Escala (10 divisões totais)
\foreach \y in {0.6, 1.2, 1.8, 2.4, 3.0, 3.6, 4.2, 4.8, 5.4, 6.0} {
    \draw[thick, marcadores] (4, \y) -- (3.6, \y);
}
% Textos com máscara fill=white para leitura limpa
\node[right, fill=white, inner sep=2pt, font=\bfseries] at (4.2, 6.0) {1 L};
\node[right, fill=white, inner sep=2pt, font=\bfseries] at (4.2, 3.0) {0,5 L};
\end{tikzpicture}
\end{adjustbox}
\end{center}

Com base na distribuição proporcional da escala visual, identifique o volume decimal de reagente contido no recipiente e converta esse valor para uma fração irredutível.
\par\vspace{\espacoQuestao}

\subsection*{Capítulo 2: Transformações e Comparações}
\vspace{-0.2cm}\noindent\rule{\linewidth}{1.5pt} 
\par\vspace{\espacoPosTitulo}
\textbf{6.} A reta numérica é o modelo geométrico perfeito para compreender as desigualdades no conjunto dos Números Racionais. Julgue as afirmativas a seguir em relação à ordenação estrutural desses números, assinalando V para verdadeiro ou F para falso.

Apresente a correção matemática para as sentenças falsas.
\begin{enumerate}[label=\Roman*.]
    \item Na reta numérica, o número \textbf{-2,54} está posicionado à direita do número \textbf{-2,5}, portanto, é maior.
    \item O número decimal \textbf{0,35} está exatamente no ponto médio entre os números \textbf{0,3} e \textbf{0,4}.
    \item Não existe nenhum número decimal exato (com duas casas decimais) compreendido entre \textbf{3,14} e \textbf{3,15}.
\end{enumerate}
\par\vspace{\espacoQuestao}

\noindent\textcolor{CorLinha}{\rule{\linewidth}{0.5pt}} 
\par\vspace{\espacoPreQuestao}
\textbf{7.} O Índice de Massa Corporal (IMC) é uma relação que pode gerar números decimais não exatos. Suponha que, para determinar um "Fator de Risco" cardíaco, um hospital faça a divisão da massa do paciente (em kg) pela sua idade (em anos). O paciente Marcos tem \textbf{87 kg} e \textbf{40} anos de idade.

Calcule o valor decimal exato do "Fator de Risco" de Marcos. Em seguida, represente como você explicaria a posição desse número na reta numérica, dizendo entre quais inteiros ele se encontra e de qual deles está mais próximo.
\par\vspace{\espacoQuestao}

\noindent\textcolor{CorLinha}{\rule{\linewidth}{0.5pt}} 
\par\vspace{\espacoPreQuestao}
\textbf{8. (Estilo VUNESP)} - Em um projeto arquitetônico corporativo, um grande salão retangular será dividido paralelamente em duas salas de reuniões menores. A planta baixa abaixo contém a distribuição estrutural da obra. A equipe de engenharia precisa conhecer a cota transversal da Sala B para instalar o forro acústico.

\begin{center}
\begin{adjustbox}{max width=\linewidth}
\begin{tikzpicture}
% --- APLICANDO DROP SHADOW E CORES HEX ---
\definecolor{parede}{HTML}{2C3E50}
\definecolor{pisoA}{HTML}{D5DBDB}
\definecolor{pisoB}{HTML}{EAEDED}

% Sombra e Fundo projetados e BEM expandidos (Evita cortes nas cotas)
\fill[black!20, rounded corners=4pt, drop shadow={opacity=0.2, shadow xshift=4pt, shadow yshift=-4pt}] (-1.5, -2.0) rectangle (10.5, 6.5);
\fill[white, rounded corners=4pt] (-1.5, -2.0) rectangle (10.5, 6.5);

% Pisos
\fill[pisoA] (0, 0) rectangle (5, 5);
\fill[pisoB] (5, 0) rectangle (8, 5);

% Paredes
\draw[line width=5pt, parede] (0,0) rectangle (8,5);
\draw[line width=5pt, parede] (5,0) -- (5,5);

% Cotas com textos interrompendo a linha (fill=white)
\draw[|<->|, thick] (0, -1.0) -- (5, -1.0) node[midway, fill=white, inner sep=3pt] {\textbf{5,45 m}};
\draw[|<->|, thick] (5, -1.0) -- (8, -1.0) node[midway, fill=white, inner sep=3pt] {\textbf{x m}};
\draw[|<->|, thick] (9.0, 0) -- (9.0, 5) node[midway, fill=white, inner sep=3pt] {\textbf{4,80 m}};

% Textos internos protegidos com fill na cor do piso
\node[font=\bfseries, fill=pisoA, inner sep=2pt] at (2.5, 2.5) {SALA A};
\node[font=\bfseries, fill=pisoB, inner sep=2pt] at (6.5, 2.5) {SALA B};
\node[font=\bfseries, text=parede, fill=pisoA, inner sep=2pt] at (2.5, 4.3) {Área Total = 41,76 m$^2$};
\end{tikzpicture}
\end{adjustbox}
\end{center}

Determine o valor da medida \textbf{x}, correspondente à largura da Sala B. Expresse a resposta final na forma de um número misto (parte inteira e fração irredutível).
\par\vspace{\espacoQuestao}

\noindent\textcolor{CorLinha}{\rule{\linewidth}{0.5pt}} 
\par\vspace{\espacoPreQuestao}
\textbf{9. (ENEM - Adaptada)} - Em uma prova de automobilismo, os tempos de volta dos pilotos são registrados com precisão de milésimos de segundo para evitar empates técnicos. Durante a volta classificatória, o sistema de telemetria registrou os seguintes dados das parciais de um setor da pista para quatro carros:

Carro W: \textbf{22,093} s\\
Carro X: \textbf{22,30} s\\
Carro Y: \textbf{22,039} s\\
Carro Z: \textbf{22,1} s

O diretor de prova solicitou que o engenheiro organizasse o grid de largada desse setor em ordem crescente de tempos (do mais rápido para o mais lento). Qual deve ser a sequência correta repassada aos pilotos?
\begin{enumerate}[label=\alph*)]
    \item Y, W, Z, X
    \item Y, Z, W, X
    \item Z, W, Y, X
    \item W, Y, Z, X
    \item X, Z, W, Y
\end{enumerate}
\par\vspace{\espacoQuestao}

\noindent\textcolor{CorLinha}{\rule{\linewidth}{0.5pt}} 
\par\vspace{\espacoPreQuestao}
\textbf{10. (Estilo FUVEST)} - No estudo analítico dos conjuntos numéricos, a relação de ordem inverte-se em determinados intervalos. Considere as variáveis algébricas $ a = 0,16 $ e $ b = 0,2 $.

Um estudante constrói as razões $ K_1 = \mfrac{a}{b} $ e $ K_2 = \mfrac{b}{a} $. Realize os cálculos necessários e responda: qual das duas grandezas ($ K_1 $ ou $ K_2 $) assumirá a posição mais à direita em uma reta numérica padrão? Justifique, apresentando a transformação dos valores decimais para suas respectivas frações irredutíveis.
\par\vspace{\espacoQuestao}

\subsection*{Capítulo 3: Operações com Decimais}
\vspace{-0.2cm}\noindent\rule{\linewidth}{1.5pt} 
\par\vspace{\espacoPosTitulo}
\textbf{11.} Uma transação bancária internacional envolve tarifas e conversões precisas. Um cliente possui um saldo de \textbf{R\$ 450,75} em sua conta corrente. Em um único dia, ele realizou as seguintes operações:
\begin{enumerate}[label=\arabic*.]
    \item Pagou uma fatura de cartão de crédito no valor de \textbf{R\$ 189,90}.
    \item Recebeu uma transferência (PIX) no valor de \textbf{R\$ 34,50}.
    \item Teve debitada uma tarifa de manutenção de conta de \textbf{R\$ 12,85}.
\end{enumerate}

Modele matematicamente a expressão numérica que representa o fluxo de caixa do cliente e calcule o saldo final exato disponível ao final do dia.
\par\vspace{\espacoQuestao}

\noindent\textcolor{CorLinha}{\rule{\linewidth}{0.5pt}} 
\par\vspace{\espacoPreQuestao}
\textbf{12.} Na auditoria das notas fiscais de um restaurante, o gerente encontrou um cupom fiscal danificado, onde a impressora falhou e ocultou dois valores importantes (identificados pelas letras \textbf{Y} e \textbf{Z}).

\begin{center}
\begin{adjustbox}{max width=\linewidth}
\begin{tikzpicture}
% --- APLICANDO DROP SHADOW E CORES HEX ---
\definecolor{papel}{HTML}{F9E79F} % Tom de papel amarelo
\definecolor{texto}{HTML}{34495E}
\definecolor{faixa}{HTML}{F5B041}

% LAYER 1: Sombra Expandida e Corpo do Cupom ajustados
\fill[black!20, rounded corners=4pt, drop shadow={opacity=0.3, shadow xshift=4pt, shadow yshift=-4pt}] (-0.5, -0.5) rectangle (7.5, 6.0);
\fill[papel, rounded corners=4pt] (-0.5, -0.5) rectangle (7.5, 6.0);
\fill[faixa, rounded corners=4pt] (-0.5, 4.8) rectangle (7.5, 6.0);

% LAYER 2: Textos protegidos com máscara de cor de fundo
\node[texto, font=\bfseries, fill=faixa, inner sep=2pt] at (3.5, 5.4) {CUPOM FISCAL ELETRÔNICO};

\node[texto, right, fill=papel, inner sep=2pt] at (0.2, 4.2) {01. Carne Bovina (2,5 kg) ........ R\$ 92,25};
\node[texto, right, fill=papel, inner sep=2pt] at (0.2, 3.4) {02. Azeite de Oliva (1 L) ........... R\$ 45,90};
\node[texto, right, fill=papel, inner sep=2pt] at (0.2, 2.6) {03. Arroz Branco (5 kg) ............ R\$ \textbf{Y}};

\draw[dashed, thick, texto] (0.2, 2.0) -- (6.8, 2.0);

\node[texto, right, font=\bfseries, fill=papel, inner sep=2pt] at (0.2, 1.2) {SUBTOTAL .............................. R\$ 168,00};
\node[texto, right, font=\bfseries, fill=papel, inner sep=2pt] at (0.2, 0.4) {TROCO (Pgto R\$ 200) ............ R\$ \textbf{Z}};
\end{tikzpicture}
\end{adjustbox}
\end{center}

O auditor financeiro precisa reconstruir o documento para o balanço mensal. Analisando as relações matemáticas de soma e subtração implícitas no documento, os valores exatos de \textbf{Y} (preço do arroz) e \textbf{Z} (valor do troco entregue ao cliente) são, respectivamente:
\begin{enumerate}[label=\alph*)]
    \item Y = 29,85 e Z = 32,00
    \item Y = 30,85 e Z = 32,00
    \item Y = 29,85 e Z = 42,00
    \item Y = 31,15 e Z = 31,00
    \item Y = 29,85 e Z = 31,00
\end{enumerate}
\par\vspace{\espacoQuestao}
\noindent\textcolor{CorLinha}{\rule{\linewidth}{0.5pt}} 
\par\vspace{\espacoPreQuestao}
\textbf{13.} Uma farmácia de manipulação adota um rigoroso controle de massa em sua linha de produção. Para um lote específico de medicação neurológica, a receita exige exatamente \textbf{0,45 g} de princípio ativo por cápsula. 

Se o laboratório recebeu uma encomenda para a produção contínua de \textbf{1000} cápsulas idênticas, determine a massa total de princípio ativo necessária. Em sua resposta, explique a regra prática de deslocamento da vírgula ao multiplicar um número decimal por potências de dez.
\par\vspace{\espacoQuestao}

\noindent\textcolor{CorLinha}{\rule{\linewidth}{0.5pt}}
\par\vspace{\espacoPreQuestao}
\textbf{14.} Durante a resolução de uma prova, o aluno fictício Carlos encontrou a operação de divisão $ \mfrac{4,8}{0,5} $. Ele escreveu a seguinte justificativa em sua folha de respostas: \textit{"Como estou dividindo por 0,5, que é exatamente a metade de 1, o resultado final deve ser a metade de 4,8, ou seja, 2,4"}.
\begin{boxresponda}
\textbf{Responda:} Identifique e explique o erro conceitual cometido por Carlos sobre a divisão de decimais. Em seguida, demonstre a manipulação correta das casas decimais para revelar o verdadeiro resultado desta operação.
\end{boxresponda}
\par\vspace{\espacoQuestao}

\noindent\textcolor{CorLinha}{\rule{\linewidth}{0.5pt}} 
\par\vspace{\espacoPreQuestao}
\textbf{15.} Um estúdio de marcenaria está produzindo prateleiras modulares. O projeto exige que uma prancha bruta de cedro seja fracionada em partes estritamente iguais. O esquema ilustra a prancha antes do corte, com sua cota total aferida a laser.

\begin{center}
\begin{adjustbox}{max width=\linewidth}
\begin{tikzpicture}
\definecolor{madeira}{HTML}{D35400}
\definecolor{fundo}{HTML}{FFFFFF}
\definecolor{cota}{HTML}{2C3E50}

% LAYER 1: Margem e Sombra
\fill[black!15, rounded corners=4pt, drop shadow={opacity=0.2, shadow xshift=3pt, shadow yshift=-3pt}] (-1, -1.5) rectangle (9, 2.0);
\fill[fundo, rounded corners=4pt] (-1, -1.5) rectangle (9, 2.0);

% LAYER 2: Prancha de Madeira
\fill[madeira!80!white, rounded corners=2pt] (0, 0) rectangle (7.65, 0.8);
\draw[thick, madeira!40!black, rounded corners=2pt] (0, 0) rectangle (7.65, 0.8);

% Textura simulada
\draw[madeira!40!black, opacity=0.3] (0.5, 0.2) -- (7.0, 0.2);
\draw[madeira!40!black, opacity=0.3] (1.0, 0.6) -- (6.5, 0.6);

% LAYER 3: Cotas
\draw[|<->|, thick, cota] (0, -0.6) -- (7.65, -0.6) node[midway, fill=fundo, inner sep=2pt, font=\bfseries] {2,55 m};

% Marcação de corte
\draw[dashed, thick, cota] (2.55, 1.2) -- (2.55, -0.2);
\draw[|<->|, thick, cota] (0, 1.0) -- (2.55, 1.0) node[midway, fill=fundo, inner sep=2pt, font=\bfseries] {0,85 m};
\end{tikzpicture}
\end{adjustbox}
\end{center}

O carpinteiro ajustou a serra para criar pedaços de \textbf{0,85 m}, sem deixar margem de perda no corte. Calcule a quantidade exata de prateleiras que ele conseguirá produzir a partir desta prancha.
\par\vspace{\espacoQuestao}
\noindent\textcolor{CorLinha}{\rule{\linewidth}{0.5pt}} 
\par\vspace{0.06cm}
\textbf{16.} A logística de transporte rodoviário exige precisão na alocação de cargas volumétricas. O esquema apresenta a montagem de um palete no interior de uma van de entregas. A altura do teto do compartimento de carga estabelece um limite estrutural rigoroso que não pode ser superado.
\par
\begin{center}
\begin{adjustbox}{max width=\linewidth}
\begin{tikzpicture}
\definecolor{palete}{HTML}{8D6E63}
\definecolor{caixaA}{HTML}{5DADE2}
\definecolor{caixaB}{HTML}{48C9B0}
\definecolor{caixaC}{HTML}{F4D03F}
\definecolor{fundo}{HTML}{F2F3F4}

% LAYER 1: Fundo e Sombra
\fill[black!15, rounded corners=4pt, drop shadow={opacity=0.2, shadow xshift=3pt, shadow yshift=-3pt}] (-1, -0.3) rectangle (6.65, 5.25);
\fill[fundo, rounded corners=4pt] (-1, -0.3) rectangle (6.65, 5.25);

% Chão
\draw[ultra thick, black!70] (-0.5, 0) -- (5.5, 0);

% LAYER 2: Palete e Caixas
\fill[palete, rounded corners=1pt] (0.5, 0) rectangle (3.5, 0.4);
\fill[caixaA, rounded corners=2pt] (0.7, 0.4) rectangle (3.3, 1.9);
\node[font=\bfseries] at (2, 1.15) {CAIXA 1};

\fill[caixaB, rounded corners=2pt] (0.8, 1.9) rectangle (3.2, 3.2);
\node[font=\bfseries] at (2, 2.55) {CAIXA 2};

\fill[caixaC, rounded corners=2pt] (0.9, 3.2) rectangle (3.1, 4.8);
\node[font=\bfseries] at (2, 4.0) {CAIXA 3};

% Teto da Van
\draw[ultra thick, red!80!black, dashed] (-0.5, 4.8) -- (5.5, 4.8);
\node[above, text=red!80!black, font=\bfseries] at (2.5, 4.8) {LIMITE DO TETO};

% LAYER 3: Cotas com máscara visual
\draw[|<->|, thick] (4.0, 0) -- (4.0, 0.4) node[midway, right, fill=fundo, inner sep=2pt, font=\bfseries] {0,15 m};
\draw[|<->|, thick] (4.0, 0.4) -- (4.0, 1.9) node[midway, right, fill=fundo, inner sep=2pt, font=\bfseries] {0,65 m};
\draw[|<->|, thick] (4.0, 1.9) -- (4.0, 3.2) node[midway, right, fill=fundo, inner sep=2pt, font=\bfseries] {0,55 m};
\draw[|<->|, thick] (4.0, 3.2) -- (4.0, 4.8) node[midway, right, fill=fundo, inner sep=2pt, font=\bfseries] {x};

\draw[|<->|, ultra thick] (5.5, 0) -- (5.5, 4.8) node[midway, right, fill=fundo, inner sep=2pt, font=\bfseries] {2,10 m};
\end{tikzpicture}
\end{adjustbox}
\end{center}

Modele a equação de soma e subtração baseada nas dimensões apresentadas e determine a altura máxima exata (\textbf{x}) que a terceira caixa pode ter para atingir milimetricamente o teto da van.
\par\vspace{\espacoQuestao}

\noindent\textcolor{CorLinha}{\rule{\linewidth}{0.5pt}} 
\par\vspace{\espacoPreQuestao}
\textbf{17. (Estilo VUNESP - Refutação de Distratores)} - Um turista brasileiro em viagem ao exterior precisa adquirir moeda estrangeira. Na casa de câmbio, a cotação do dólar turismo para compra estava fixada em \textbf{R\$ 4,95}. Além do valor convertido sobre a moeda, a corretora cobra uma taxa fixa operacional (IOF) de \textbf{R\$ 15,50} por transação, independentemente do montante trocado.

O turista solicitou a conversão exata de \textbf{120} dólares. Avalie as alternativas abaixo referentes ao custo total da operação.
\begin{enumerate}[label=\alph*)]
    \item R\$ 594,00
    \item R\$ 609,50
    \item R\$ 615,00
    \item R\$ 578,50
    \item R\$ 620,50
\end{enumerate}
\begin{boxresponda}
\textbf{Responda:} Qual é a alternativa correta? Comprove com cálculos por que a alternativa \textbf{a)} é um "distrator" intencional, explicando qual erro lógico de leitura do enunciado um candidato cometeria para chegar a esse valor.
\end{boxresponda}
\par\vspace{0.0cm}

\noindent\textcolor{CorLinha}{\rule{\linewidth}{0.5pt}} 
\par\vspace{\espacoPreQuestao}
\textbf{18.} A eficiência energética de um automóvel é medida pela razão entre a distância percorrida e o volume de combustível consumido. Um modelo sedã possui um tanque de combustível com capacidade máxima de \textbf{45,5} litros. Em testes de estrada, a montadora certificou que o veículo possui um rendimento médio de \textbf{12,4} km por litro.
\begin{boxdica}
Atenção à regra de deslocamento da vírgula! O número de casas decimais do produto final é igual à soma das casas decimais dos fatores multiplicados.
\end{boxdica}
Demonstre a armação da multiplicação e determine a distância máxima, em quilômetros, que este carro consegue percorrer com o tanque completamente cheio.
\par\vspace{0.0cm}

\noindent\textcolor{CorLinha}{\rule{\linewidth}{0.5pt}} 
\par\vspace{0.05cm}
\textbf{19. (Estilo FUVEST)} - A análise de escoamento de fluidos em engenharia ambiental utiliza gráficos de barras para registrar o índice pluviométrico. O gráfico apresenta a precipitação em milímetros registrada em uma microbacia durante quatro dias consecutivos.
\begin{center}
\begin{adjustbox}{max width=\linewidth}
\begin{tikzpicture}
\definecolor{fundo}{HTML}{FFFFFF}
\definecolor{barra}{HTML}{3498DB}
\definecolor{eixo}{HTML}{2C3E50}

% Fundo Seguro
\fill[black!10, rounded corners=4pt, drop shadow={opacity=0.2, shadow xshift=3pt, shadow yshift=-3pt}] (-1, -0.7) rectangle (8, 4.3);
\fill[fundo, rounded corners=4pt] (-1, -0.7) rectangle (8, 4.3);

\begin{axis}[
    width=8cm, height=5.5cm,
    axis lines=left,
    axis line style={thick, eixo},
    ymin=0, ymax=15,
    xmin=0, xmax=5,
    xtick={1,2,3,4},
    xticklabels={Dia 1, Dia 2, Dia 3, Dia 4},
    ytick={0, 2.5, 5, 7.5, 10, 12.5},
    ymajorgrids=true,
    grid style={dashed, gray!50},
    bar width=20pt,
    enlarge x limits=0.15,
    nodes near coords,
    every node near coord/.append style={fill=fundo, inner sep=2pt, font=\bfseries},
    ylabel={Precipitação (mm)},
    ylabel style={font=\bfseries, fill=fundo, inner sep=2pt},
    tick label style={font=\bfseries}
]

\addplot[ybar, fill=barra!80!white, draw=barra!50!black, thick, rounded corners=1pt] coordinates {
    (1, 12.4)
    (2, 8.5)
    (3, 10.25)
    (4, 5.8)
};
\end{axis}
\end{tikzpicture}
\end{adjustbox}
\end{center}
\par
Com base no comportamento pluviométrico diário extraído do gráfico, determine a média aritmética de chuva (em mm) que incidiu sobre a microbacia nestes quatro dias.
\par\vspace{\espacoQuestao}

\noindent\textcolor{CorLinha}{\rule{\linewidth}{0.5pt}}
\par\vspace{\espacoPreQuestao}
\textbf{20. (ENEM - Adaptada)} - No desenvolvimento de uma liga metálica aeroespacial, um engenheiro químico mistura \textbf{2,5 kg} do Metal Alfa, que custa \textbf{R\$ 140,40} o quilograma, com \textbf{1,5 kg} do Metal Beta, cujo custo é \textbf{R\$ 180,60} o quilograma. Após o derretimento, forma-se um bloco homogêneo de \textbf{4,0 kg}.

Ao tabelar o preço da nova liga, um estagiário afirmou que o preço do quilograma final seria de \textbf{R\$ 160,50}, pois calculou a média simples dos preços: $ \mfrac{140,40 + 180,60}{2} $.
\begin{boxresponda}
\textbf{Responda:} O raciocínio do estagiário está estruturalmente equivocado para esta modelagem. Refute o método utilizado por ele e apresente o cálculo completo que determina o verdadeiro valor de custo de um quilograma da nova liga.
\end{boxresponda}
\par\vspace{\espacoQuestao}

\subsection*{Capítulo 4: Porcentagem}
\vspace{-0.2cm}\noindent\rule{\linewidth}{1.5pt} 
\par\vspace{\espacoPosTitulo}
\textbf{21.} A representação geométrica através de malhas é uma das formas mais visuais de se compreender a essência da porcentagem, que significa matematicamente "uma fração centesimal". O diagrama ilustra um piso de \textbf{100} lajotas, onde uma parte exata foi revestida com resina especial antiderrapante.

\begin{center}
\begin{adjustbox}{max width=\linewidth}
\begin{tikzpicture}
\definecolor{fundo}{HTML}{F2F4F4}
\definecolor{quadrado}{HTML}{16A085}

% Sombra e Fundo Seguro
\fill[black!15, rounded corners=4pt, drop shadow={opacity=0.2, shadow xshift=3pt, shadow yshift=-3pt}] (-1, -1) rectangle (6, 6);
\fill[fundo, rounded corners=4pt] (-1, -1) rectangle (6, 6);

% Preenchimento da área resinada (37 quadrados de 0.5x0.5)
\fill[quadrado] (0,0) rectangle (5, 1.5);
\fill[quadrado] (0,1.5) rectangle (3.5, 2.0);

% Malha 10x10
\draw[step=0.5, black!60, thick] (0,0) grid (5,5);
\draw[ultra thick, black!80] (0,0) rectangle (5,5);
\end{tikzpicture}
\end{adjustbox}
\end{center}

Com base na proporção da área resinada (em verde) em relação à área total do piso, expresse essa quantidade extraída visualmente em três formatos matemáticos distintos: na forma de taxa percentual (\%), na forma de fração irredutível e na forma de número decimal.
\par\vspace{\espacoQuestao}

\noindent\textcolor{CorLinha}{\rule{\linewidth}{0.5pt}} 
\par\vspace{\espacoPreQuestao}
\textbf{22.} O domínio das transformações entre taxas percentuais, decimais e frações é a base da alfabetização financeira moderna. Analise as afirmativas a seguir sobre essas equivalências e julgue-as assinalando V para verdadeiro ou F para falso. 

\begin{enumerate}[label=\Roman*.]
    \item A taxa percentual de \textbf{5\%} é matematicamente equivalente à fração irredutível $ \mfrac{1}{20} $.
    \item O número decimal \textbf{0,08} representa exatamente \textbf{80\%} de uma quantidade inteira.
    \item Calcular \textbf{150\%} de um determinado valor \textbf{x} é o mesmo que somar a sua própria metade ao valor original.
\end{enumerate}
Para cada afirmativa considerada falsa, apresente a conversão e o cálculo correto.
\par\vspace{\espacoQuestao}

\noindent\textcolor{CorLinha}{\rule{\linewidth}{0.5pt}}
\par\vspace{\espacoPreQuestao}
\textbf{23.} A matemática dos reajustes percentuais sucessivos costuma gerar confusões lógicas devido à alteração constante da base de cálculo. Ao observar a vitrine de uma loja, a aluna Mariana notou que uma jaqueta sofreu um aumento de \textbf{20\%} no mês de maio e, em junho, entrou em promoção com um desconto de \textbf{20\%}.
Mariana concluiu: \textit{"Se o preço subiu 20\% e depois desceu 20\%, a jaqueta voltou a custar exatamente o que custava antes do aumento"}.
\begin{boxresponda}
\textbf{Responda:} A afirmação de Mariana está correta ou incorreta? Prove algebricamente ou através de um exemplo numérico (supondo que a jaqueta custasse R\$ 100,00) qual foi o impacto real no preço final após as duas alterações sucessivas.
\end{boxresponda}
\par\vspace{\espacoQuestao}

\noindent\textcolor{CorLinha}{\rule{\linewidth}{0.5pt}} 
\par\vspace{\espacoPreQuestao}
\textbf{24.} A capacidade de armazenamento de energia das baterias de dispositivos móveis é medida em miliampere-hora (mAh). O painel de controle de um smartphone exibe o nível atual. Sabe-se que a bateria deste aparelho possui uma capacidade total de \textbf{4.500 mAh}.

\begin{center}
\begin{adjustbox}{max width=\linewidth}
\begin{tikzpicture}
\definecolor{fundo}{HTML}{FFFFFF}
\definecolor{batbg}{HTML}{D5DBDB}
\definecolor{batfill}{HTML}{F39C12}

% Fundo e Sombra Expandidos
\fill[black!15, rounded corners=4pt, drop shadow={opacity=0.2, shadow xshift=3pt, shadow yshift=-3pt}] (-1.5, -1.5) rectangle (7.5, 2);
\fill[fundo, rounded corners=4pt] (-1.5, -1.5) rectangle (7.5, 2);

% Corpo da bateria
\draw[ultra thick, black!80, rounded corners=3pt] (0, -0.5) rectangle (5, 1);
\fill[black!80, rounded corners=1pt] (5.1, -0.1) rectangle (5.3, 0.6);
\fill[batbg, rounded corners=2pt] (0.1, -0.4) rectangle (4.9, 0.9);

% Carga visual (68% de 4.8)
\fill[batfill, rounded corners=2pt] (0.1, -0.4) rectangle (3.36, 0.9);

\node[font=\bfseries, text=black!90] at (2.5, 0.25) {68\% de Carga};
\end{tikzpicture}
\end{adjustbox}
\end{center}

Lendo os dados visuais do painel, determine a quantidade de carga exata, em mAh, que \textbf{já foi consumida} (ou seja, que foi gasta pelo sistema) desde o momento em que o aparelho foi tirado da tomada com 100\%.
\par\vspace{\espacoQuestao}

\noindent\textcolor{CorLinha}{\rule{\linewidth}{0.5pt}} 
\par\vspace{\espacoPreQuestao}
\textbf{25. (Estilo FUVEST - Refutação)} - Técnicas de marketing no varejo frequentemente mascaram descontos percentuais reais usando promoções em unidades físicas. Um supermercado lançou a famosa promoção para uma marca de cervejas artesanais: \textbf{"LEVE 5, PAGUE 4"}. 

O cliente desatento pode acreditar, erroneamente, que essa promoção resulta em uma redução direta de 25\% no ato da compra.
\begin{enumerate}[label=\alph*)]
    \item 20\%
    \item 25\%
    \item 15\%
    \item 10\%
    \item 30\%
\end{enumerate}
\begin{boxresponda}
\textbf{Responda:} Qual é a alternativa que apresenta a taxa percentual exata do desconto real obtido pelo consumidor? Por que o raciocínio mental que leva aos 25\% está equivocado em relação ao conceito de base de cálculo?
\end{boxresponda}
\par\vspace{\espacoQuestao}
\noindent\textcolor{CorLinha}{\rule{\linewidth}{0.5pt}} 
\par\vspace{\espacoPreQuestao}
\textbf{26.} O planejamento orçamentário familiar exige uma distribuição rigorosa das despesas. O gráfico de setores (pizza) representa o destino da renda mensal de um casal, cujo rendimento total líquido soma \textbf{R\$ 6.000,00}.

\begin{center}
\begin{adjustbox}{max width=\linewidth}
\begin{tikzpicture}
% --- CORES CUSTOMIZADAS TIKZ ---
\definecolor{fundo}{HTML}{F8F9F9}
\definecolor{setor1}{HTML}{3498DB} % Moradia 35%
\definecolor{setor2}{HTML}{E74C3C} % Alimentação x% -> será 30%
\definecolor{setor3}{HTML}{F1C40F} % Transporte 15%
\definecolor{setor4}{HTML}{2ECC71} % Educação 20%

% Sombra e Fundo
\fill[black!15, rounded corners=4pt, drop shadow={opacity=0.2, shadow xshift=3pt, shadow yshift=-3pt}] (-3.5, -3.2) rectangle (4.5, 3.2);
\fill[fundo, rounded corners=4pt] (-3.5, -3.2) rectangle (4.5, 3.2);

% Setores (Raio 2)
% Moradia 35% = 126 graus (0 a 126)
\draw[fill=setor1, draw=white, thick] (0,0) -- (0:2.2) arc (0:126:2.2) -- cycle;
% Educação 20% = 72 graus (126 a 198)
\draw[fill=setor4, draw=white, thick] (0,0) -- (126:2.2) arc (126:198:2.2) -- cycle;
% Transporte 15% = 54 graus (198 a 252)
\draw[fill=setor3, draw=white, thick] (0,0) -- (198:2.2) arc (198:252:2.2) -- cycle;
% Alimentação x% = 30% = 108 graus (252 a 360)
\draw[fill=setor2, draw=white, thick] (0,0) -- (252:2.2) arc (252:360:2.2) -- cycle;

% Legendas com máscara de fundo (fill=fundo)
\node[fill=fundo, inner sep=2pt, font=\bfseries] at (63:2.7) {Moradia: 35\%};
\node[fill=fundo, inner sep=2pt, font=\bfseries] at (162:2.8) {Educação: 20\%};
\node[fill=fundo, inner sep=2pt, font=\bfseries] at (225:2.8) {Transporte: 15\%};
\node[fill=fundo, inner sep=2pt, font=\bfseries] at (306:2.8) {Alimentação: x\%};
\end{tikzpicture}
\end{adjustbox}
\end{center}

Com base nas informações visuais, calcule primeiro a taxa percentual destinada à Alimentação (\textbf{x\%}). Em seguida, determine o valor financeiro absoluto, em reais, que essa família gasta exclusivamente com esse setor.
\par\vspace{\espacoQuestao}

\noindent\textcolor{CorLinha}{\rule{\linewidth}{0.5pt}} 
\par\vspace{\espacoPreQuestao}
\textbf{27.} Em uma escola de idiomas de grande porte, estão matriculados \textbf{1.200} alunos nos mais diversos níveis. O sistema acadêmico aponta que \textbf{35\%} desse total cursam o idioma Espanhol. Dentre os alunos que cursam Espanhol, sabe-se que \textbf{20\%} estão matriculados nas turmas do nível avançado.

Determine o número exato de alunos da escola que estudam Espanhol no nível avançado. Demonstre as duas etapas do seu cálculo.
\par\vspace{\espacoQuestao}

\noindent\textcolor{CorLinha}{\rule{\linewidth}{0.5pt}} 
\par\vspace{\espacoPreQuestao}
\textbf{28.} O professor de matemática pediu para a turma calcular o valor final de um tênis que custava \textbf{R\$ 250,00} e sofreu um acréscimo de \textbf{15\%}. O aluno fictício João montou a seguinte operação em seu caderno: \textit{"Como 15\% é igual a 0,15, basta eu somar os valores absolutos de forma direta. O novo preço será $ 250,00 + 0,15 = \text{R\$ } 250,15 $."}
\begin{boxresponda}
\textbf{Responda:} Qual erro conceitual e estrutural João cometeu ao interpretar a porcentagem como um número decimal absoluto na soma? Corrija o cálculo e encontre o verdadeiro novo preço do tênis, explicando a função da "Base de Cálculo".
\end{boxresponda}
\par\vspace{\espacoQuestao}

\noindent\textcolor{CorLinha}{\rule{\linewidth}{0.5pt}} 
\par\vspace{\espacoPreQuestao}
\textbf{29. (ENEM - Adaptada)} - A inflação causa uma remarcação natural nos preços das prateleiras de um supermercado. No início de um semestre, um pacote de café de marca popular custava \textbf{R\$ 16,00}. Ao final do mesmo semestre, sem alteração no tamanho da embalagem, o mesmo produto passou a ser etiquetado por \textbf{R\$ 18,40}.

Qual foi a taxa percentual de aumento que incidiu sobre o preço original deste produto?
\begin{enumerate}[label=\alph*)]
    \item 12\%
    \item 15\%
    \item 18\%
    \item 24\%
    \item 240\%
\end{enumerate}
\begin{boxresponda}
\textbf{Responda:} Assinale a alternativa correta demonstrando o seu cálculo. Em seguida, prove matematicamente por que a alternativa \textbf{e)} é um "distrator lógico" e explique que tipo de divisão invertida um aluno desatento faria para chegar a esse valor absurdo.
\end{boxresponda}
\par\vspace{\espacoQuestao}

\noindent\textcolor{CorLinha}{\rule{\linewidth}{0.5pt}} 
\par\vspace{\espacoPreQuestao}
\textbf{30.} A interface de um sistema operacional exibe uma barra de progresso durante o download de uma atualização crítica de software. O pacote total de arquivos do sistema pesa \textbf{18 GB} (Gigabytes).

\begin{center}
\begin{adjustbox}{max width=\linewidth}
\begin{tikzpicture}
\definecolor{fundo}{HTML}{FFFFFF}
\definecolor{barrabg}{HTML}{E5E8E8}
\definecolor{barrafill}{HTML}{8E44AD}

% Fundo e Sombra
\fill[black!15, rounded corners=4pt, drop shadow={opacity=0.2, shadow xshift=3pt, shadow yshift=-3pt}] (-1, -1.5) rectangle (9, 2);
\fill[fundo, rounded corners=4pt] (-1, -1.5) rectangle (9, 2);

% Texto Superior
\node[anchor=west, font=\bfseries] at (0, 1.2) {Baixando Atualização: Total de 18 GB};

% Barra de fundo
\fill[barrabg, rounded corners=4pt] (0, 0) rectangle (8, 0.6);
% Preenchimento (45% de 8 = 3.6)
\fill[barrafill, rounded corners=4pt] (0, 0) rectangle (3.6, 0.6);

% Texto dentro/sobre a barra (fill=white para contraste)
\node[fill=barrafill, text=white, inner sep=2pt, font=\bfseries] at (1.8, 0.3) {45\%};
\end{tikzpicture}
\end{adjustbox}
\end{center}

Neste exato momento da transferência, quantos Gigabytes \textbf{ainda faltam} ser baixados para que a instalação do sistema seja totalmente concluída?
\par\vspace{\espacoQuestao}

\noindent\textcolor{CorLinha}{\rule{\linewidth}{0.5pt}} 
\par\vspace{\espacoPreQuestao}
\textbf{31.} Na construção civil, alterações nos projetos geram impactos compostos na área final. A planta de um terreno retangular sofreu uma revisão municipal: a medida do seu comprimento foi \textbf{aumentada em 20\%}, enquanto a medida da sua largura foi obrigatoriamente \textbf{reduzida em 10\%}. 

Considerando as duas alterações simultâneas, modele algebricamente o problema (ou adote medidas fictícias para teste) e demonstre o que ocorrerá com a área final do terreno em relação à área inicial. Ela aumentou ou diminuiu? De qual percentual?
\par\vspace{\espacoQuestao}

\noindent\textcolor{CorLinha}{\rule{\linewidth}{0.5pt}} 
\par\vspace{\espacoPreQuestao}
\textbf{32.} Em uma vitrine de calçados, um par de tênis para corrida está sendo anunciado com uma etiqueta promocional marcando \textbf{R\$ 204,00}. O vendedor informa que esse valor já embute um desconto absoluto de \textbf{15\%} em relação ao preço original. 

O estudante Lucas tenta descobrir o preço original e faz a seguinte conta: \textit{"Se a loja deu desconto de 15\%, basta eu calcular 15\% de R\$ 204,00 e somar de volta. Assim: $ 204 \times 0,15 = 30,60 $. Portanto, o preço original era $ 204 + 30,60 = 234,60 $."}
\begin{boxresponda}
\textbf{Responda:} Lucas cometeu um dos maiores erros da matemática financeira. Explique detalhadamente por que a "Base de Cálculo" utilizada por ele está errada. Em seguida, aplique a modelagem correta e determine o verdadeiro preço original do tênis.
\end{boxresponda}
\par\vspace{\espacoQuestao}

\noindent\textcolor{CorLinha}{\rule{\linewidth}{0.5pt}} 
\par\vspace{\espacoPreQuestao}
\textbf{33.} No mercado financeiro, a volatilidade de um ativo exige atenção aos percentuais diários de alta e queda. O gráfico de linha abaixo descreve a flutuação do preço de uma cota de fundo imobiliário ao longo de quatro dias úteis.

\begin{center}
\begin{adjustbox}{max width=\linewidth}
\begin{tikzpicture}
\definecolor{fundo}{HTML}{FDFEFE}
\definecolor{linha}{HTML}{E67E22}
\definecolor{eixo}{HTML}{2C3E50}

% Fundo e Sombra Seguros
\fill[black!15, rounded corners=4pt, drop shadow={opacity=0.2, shadow xshift=3pt, shadow yshift=-3pt}] (-1.5, -1.5) rectangle (7.5, 6.0);
\fill[fundo, rounded corners=4pt] (-1.5, -1.5) rectangle (7.5, 6.0);

\begin{axis}[
    width=8cm, height=6cm,
    axis lines=left,
    axis line style={ultra thick, eixo},
    ymin=100, ymax=160,
    xmin=0.5, xmax=4.5,
    xtick={1,2,3,4},
    xticklabels={Seg, Ter, Qua, Qui},
    ytick={100, 115, 125, 140, 150},
    ymajorgrids=true,
    grid style={dashed, gray!50},
    enlarge x limits=0.05,
    ylabel={Preço da Cota (R\$)},
    ylabel style={font=\bfseries, fill=fundo, inner sep=2pt},
    tick label style={font=\bfseries}
]

% Linha principal
\addplot[color=linha, ultra thick, mark=*, mark size=3pt, mark options={fill=white, draw=linha, thick}] coordinates {
    (1, 125)
    (2, 150)
    (3, 114)
    (4, 140)
};

% Rótulos nos pontos com máscara de fundo
\node[fill=fundo, inner sep=2pt, font=\bfseries] at (axis cs:1, 131) {R\$ 125};
\node[fill=fundo, inner sep=2pt, font=\bfseries] at (axis cs:2, 156) {R\$ 150};
\node[fill=fundo, inner sep=2pt, font=\bfseries] at (axis cs:3, 108) {R\$ 114};
\node[fill=fundo, inner sep=2pt, font=\bfseries] at (axis cs:4, 146) {R\$ 140};
\end{axis}
\end{tikzpicture}
\end{adjustbox}
\end{center}

Analisando a trajetória da cota, determine a taxa percentual de desvalorização (queda) sofrida exclusivamente na passagem da terça-feira para a quarta-feira.
\par\vspace{\espacoQuestao}

\noindent\textcolor{CorLinha}{\rule{\linewidth}{0.5pt}} 
\par\vspace{\espacoPreQuestao}
\textbf{34.} Analise as afirmativas a seguir sobre a matemática dos acréscimos e descontos. Escreva V para verdadeiro ou F para falso.
\begin{enumerate}[label=\Roman*.]
    \item Aumentar o preço de uma mercadoria em \textbf{100\%} significa que o novo preço será o dobro do preço original.
    \item Aplicar um desconto de \textbf{10\%} e, logo em seguida, um acréscimo de \textbf{10\%} sobre um produto de \textbf{R\$ 200,00} fará com que o preço final retorne aos exatos \textbf{R\$ 200,00}.
    \item Multiplicar qualquer número decimal por \textbf{1,25} é o equivalente algébrico direto a aplicar um acréscimo de \textbf{25\%} sobre este número.
\end{enumerate}
Apresente a correção matemática completa ou um contraexemplo numérico que justifique a afirmativa que você julgou ser falsa.
\par\vspace{\espacoQuestao}

\noindent\textcolor{CorLinha}{\rule{\linewidth}{0.5pt}} 
\par\vspace{\espacoPreQuestao}
\textbf{35. (Estilo VUNESP)} - Em uma indústria de cosméticos, um perfumista precisa preparar um lote de loção adstringente. Ele decide misturar \textbf{40} litros de uma solução que contém \textbf{30\%} de extrato de aloe vera com \textbf{60} litros de uma segunda solução que contém \textbf{50\%} do mesmo extrato. 

Após a mistura completa dessas duas soluções em um único reservatório, qual será a porcentagem exata de extrato de aloe vera no volume final obtido?
\begin{enumerate}[label=\alph*)]
    \item 40\%
    \item 42\%
    \item 45\%
    \item 38\%
    \item 80\%
\end{enumerate}
\begin{boxresponda}
\textbf{Responda:} Assinale a alternativa correta através do seu cálculo. Por que um candidato que assinalasse a alternativa \textbf{a)} estaria cometendo um grave erro de modelagem química e matemática? Explique.
\end{boxresponda}
\par\vspace{\espacoQuestao}

\noindent\textcolor{CorLinha}{\rule{\linewidth}{0.5pt}} 
\par\vspace{\espacoPreQuestao}
\textbf{36.} A diferença conceitual entre margem de lucro sobre o custo e margem de lucro sobre a venda é uma das métricas mais vitais da contabilidade. Uma fábrica produz camisetas esportivas a um custo de produção de \textbf{R\$ 40,00} por unidade. Essas mesmas camisetas são repassadas para as lojas e vendidas ao consumidor final por \textbf{R\$ 50,00}.
\begin{boxdica}
A margem "sobre o custo" utiliza o valor de fabricação como referencial ($ 100\% $). Já a margem "sobre a venda" utiliza o preço de vitrine como referencial principal de base.
\end{boxdica}
Calcule e responda separadamente: qual é a taxa percentual de lucro obtida em relação ao preço de custo? E qual é a taxa percentual de lucro considerando-se o preço de venda?
\par\vspace{\espacoQuestao}

\noindent\textcolor{CorLinha}{\rule{\linewidth}{0.5pt}} 
\par\vspace{\espacoPreQuestao}
\textbf{37.} Na pesquisa de satisfação de um refeitório corporativo, diagramas lógicos são utilizados para mapear a aceitação do cardápio pelos funcionários. O diagrama de Venn exibe o perfil percentual de preferência entre duas opções de refeição.

\begin{center}
\begin{adjustbox}{max width=\linewidth}
\begin{tikzpicture}
\definecolor{fundo}{HTML}{F4F6F7}
\definecolor{circA}{HTML}{E74C3C} % Carne
\definecolor{circB}{HTML}{3498DB} % Frango
\definecolor{texto}{HTML}{2C3E50}

% Sombra Projetada
\fill[black!15, rounded corners=4pt, drop shadow={opacity=0.2, shadow xshift=3pt, shadow yshift=-3pt}] (-3, -3) rectangle (7, 4);
\fill[fundo, rounded corners=4pt] (-3, -3) rectangle (7, 4);
\draw[thick, black!70, rounded corners=4pt] (-3, -3) rectangle (7, 4);

% Universo
\node[anchor=north west, font=\bfseries] at (-2.8, 3.8) {Total Entrevistado: 100\%};
\node[anchor=south east, font=\bfseries, fill=fundo, inner sep=2pt] at (6.8, -2.8) {Nenhum dos dois: 15\%};

% Círculos com opacidade
\fill[circA, opacity=0.3] (0, 0) circle (2.2);
\draw[thick, circA] (0, 0) circle (2.2);

\fill[circB, opacity=0.3] (3, 0) circle (2.2);
\draw[thick, circB] (3, 0) circle (2.2);

% Rótulos com fill=white
\node[font=\bfseries, text=texto, fill=white, inner sep=2pt, rounded corners=1pt] at (0, 2.5) {Carne (60\%)};
\node[font=\bfseries, text=texto, fill=white, inner sep=2pt, rounded corners=1pt] at (3, 2.5) {Frango (55\%)};

% Valor na intersecção
\node[font=\bfseries, text=black, fill=white, inner sep=2pt, rounded corners=1pt] at (1.5, 0) {x\%};
\end{tikzpicture}
\end{adjustbox}
\end{center}

Baseando-se nos dados apresentados, equacione o problema da intersecção de conjuntos e determine a taxa percentual de funcionários (\textbf{x\%}) que declararam gostar simultaneamente das duas opções.
\par\vspace{\espacoQuestao}

\noindent\textcolor{CorLinha}{\rule{\linewidth}{0.5pt}} 
\par\vspace{\espacoPreQuestao}
\textbf{38. (Desafio ITA/FUVEST)} - A conservação de alimentos por desidratação (secagem) altera drasticamente a concentração de massa dos compostos. Um produtor colheu \textbf{100 kg} de cogumelos frescos. Análises laboratoriais mostram que \textbf{99\%} da massa desses cogumelos frescos é composta apenas por água.

Após algumas horas em uma estufa de desidratação contínua, os cogumelos perderam parte dessa umidade por evaporação, de modo que a água passou a representar exatamente \textbf{98\%} da massa total dos cogumelos, agora parcialmente secos. 

Sabendo que apenas a água evaporou e a "massa sólida" do cogumelo permaneceu intacta, calcule a massa total exata (em kg) desse lote de cogumelos logo após a retirada da estufa.
\par\vspace{\espacoQuestao}
\noindent\textcolor{CorLinha}{\rule{\linewidth}{0.5pt}} 
\par\vspace{\espacoPreQuestao}
\textbf{39.} A análise de desempenho de vendas de uma startup de tecnologia é feita trimestralmente. O gráfico de colunas detalha a receita bruta (em milhares de reais) gerada pelo principal aplicativo da empresa durante os três primeiros meses do ano.

\begin{center}
\begin{adjustbox}{max width=\linewidth}
\begin{tikzpicture}
\definecolor{fundo}{HTML}{FFFFFF}
\definecolor{barra}{HTML}{9B59B6}
\definecolor{eixo}{HTML}{34495E}

% Fundo e Sombra Seguros
\fill[black!15, rounded corners=4pt, drop shadow={opacity=0.2, shadow xshift=3pt, shadow yshift=-3pt}] (-1.5, -1.5) rectangle (7.5, 5.5);
\fill[fundo, rounded corners=4pt] (-1.5, -1.5) rectangle (7.5, 5.5);

\begin{axis}[
    width=7cm, height=5.5cm,
    axis lines=left,
    axis line style={ultra thick, eixo},
    ymin=0, ymax=150,
    xmin=0, xmax=4,
    xtick={1,2,3},
    xticklabels={Jan, Fev, Mar},
    ytick={0, 50, 100, 150},
    ymajorgrids=true,
    grid style={dashed, gray!40},
    bar width=25pt,
    enlarge x limits=0.2,
    nodes near coords,
    every node near coord/.append style={fill=white, inner sep=2pt, font=\bfseries},
    ylabel={Receita (Milhares de R\$)},
    ylabel style={font=\bfseries, fill=white, inner sep=2pt},
    tick label style={font=\bfseries}
]

\addplot[ybar, fill=barra!80!white, draw=barra, ultra thick, rounded corners=2pt] coordinates {
    (1, 80)
    (2, 100)
    (3, 125)
};
\end{axis}
\end{tikzpicture}
\end{adjustbox}
\end{center}

Determine a taxa percentual de crescimento da receita de fevereiro para março. Em seguida, explique se o aumento percentual de janeiro para fevereiro foi igual, maior ou menor que o aumento percentual de fevereiro para março.
\par\vspace{\espacoQuestao}

\noindent\textcolor{CorLinha}{\rule{\linewidth}{0.5pt}} 
\par\vspace{\espacoPreQuestao}
\textbf{40.} O mercado automotivo utiliza tabelas de depreciação para calcular a desvalorização dos veículos ao longo do tempo. Um carro zero quilômetro foi comprado por \textbf{R\$ 90.000,00}. No primeiro ano de uso, ele sofreu uma desvalorização de \textbf{10\%}. No segundo ano, sobre o valor já depreciado, o carro perdeu mais \textbf{15\%} de seu valor de mercado.

Calcule o valor de mercado deste carro ao final do segundo ano. O desconto total acumulado nesses dois anos equivale a uma taxa percentual única de desvalorização de quantos por cento em relação ao valor original de fábrica?
\par\vspace{\espacoQuestao}

\noindent\textcolor{CorLinha}{\rule{\linewidth}{0.5pt}} 
\par\vspace{\espacoPreQuestao}
\textbf{41.} O estudante Felipe observou que a conta de energia elétrica de sua casa aumentou \textbf{10\%} em um mês e, devido a uma nova bandeira tarifária, aumentou novamente \textbf{10\%} no mês seguinte. 

Ao conversar com seu pai, Felipe afirmou: \textit{"A conta de luz sofreu um aumento total exato de 20\% em relação ao que pagávamos há dois meses, pois $ 10\% + 10\% = 20\% $."}
\begin{boxresponda}
\textbf{Responda:} Felipe cometeu o erro mais comum da matemática financeira: a soma direta de taxas relativas. Utilizando um valor inicial fictício de \textbf{R\$ 100,00} para a conta de luz, prove que a afirmação de Felipe é falsa e determine a verdadeira taxa percentual de aumento acumulado no período.
\end{boxresponda}
\par\vspace{\espacoQuestao}

\noindent\textcolor{CorLinha}{\rule{\linewidth}{0.5pt}} 
\par\vspace{\espacoPreQuestao}
\textbf{42.} No planejamento urbano e paisagístico, os municípios exigem que uma taxa mínima da área dos terrenos seja mantida como solo permeável (área verde). O projeto da residência da família Silveira foi desenhado sobre um lote retangular, com a área construída e a piscina devidamente alocadas, deixando o restante do terreno para o gramado.

\begin{center}
\begin{adjustbox}{max width=\linewidth}
\begin{tikzpicture}
\definecolor{grama}{HTML}{27AE60}
\definecolor{casa}{HTML}{95A5A6}
\definecolor{piscina}{HTML}{3498DB}
\definecolor{fundo}{HTML}{FFFFFF}

% Sombra Projetada
\fill[black!20, rounded corners=4pt, drop shadow={opacity=0.2, shadow xshift=4pt, shadow yshift=-4pt}] (-1.5, -1.5) rectangle (9.5, 6.5);
\fill[fundo, rounded corners=4pt] (-1.5, -1.5) rectangle (9.5, 6.5);

% Terreno Total (Grama no fundo)
\fill[grama!80!white, rounded corners=2pt] (0, 0) rectangle (8, 5);
\draw[ultra thick, black!80] (0, 0) rectangle (8, 5);

% Casa
\fill[casa, rounded corners=1pt] (0.5, 2.5) rectangle (4.5, 4.5);
\draw[thick, black!90] (0.5, 2.5) rectangle (4.5, 4.5);
\node[font=\bfseries, text=white] at (2.5, 3.5) {CASA};

% Piscina
\fill[piscina, rounded corners=3pt] (5.1, 0.5) rectangle (7.1, 2.5);
\draw[thick, black!90] (5.1, 0.5) rectangle (7.1, 2.5);
\node[font=\bfseries, text=white] at (6.1, 1.5) {PISCINA};

% Cotas Terreno
\draw[|<->|, thick] (0, -0.6) -- (8, -0.6) node[midway, fill=fundo, inner sep=2pt, font=\bfseries] {40 m};
\draw[|<->|, thick] (-0.6, 0) -- (-0.6, 5) node[midway, fill=fundo, inner sep=2pt, font=\bfseries] {25 m};

% Cotas Casa
\draw[|<->|, thick] (0.5, 4.8) -- (4.5, 4.8) node[midway, fill=fundo, inner sep=2pt, font=\bfseries] {20 m};
\draw[|<->|, thick] (4.8, 2.5) -- (4.8, 4.5) node[midway, fill=fundo, inner sep=2pt, font=\bfseries] {10 m};

% Cotas Piscina
\draw[|<->|, thick] (5.1, -0.2) -- (7.1, -0.2) node[midway, fill=fundo, inner sep=2pt, font=\bfseries] {10 m};
\draw[|<->|, thick] (7.5, 0.5) -- (7.5, 2.5) node[midway, fill=fundo, inner sep=2pt, font=\bfseries] {10 m};
\end{tikzpicture}
\end{adjustbox}
\end{center}

Determine a área total do terreno em metros quadrados. Em seguida, calcule a taxa percentual exata que a área verde (gramado) representa em relação à área total do lote.
\par\vspace{\espacoQuestao}

\noindent\textcolor{CorLinha}{\rule{\linewidth}{0.5pt}} 
\par\vspace{\espacoPreQuestao}
\textbf{43.} O poder de compra de um cidadão é afetado diretamente pela relação entre o reajuste salarial e a inflação do período. Em janeiro, o custo de vida mensal (despesas essenciais) da senhora Helena era de \textbf{R\$ 3.000,00}, exatamente o mesmo valor de seu salário líquido da época. 
\begin{boxdica}
A inflação age como um "acréscimo sucessivo" invisível no valor dos produtos e contas do dia a dia.
\end{boxdica}
No ano seguinte, seu salário teve um reajuste de \textbf{8\%}. Contudo, os itens que compõem suas despesas essenciais sofreram uma inflação acumulada de \textbf{11\%}. Após essas atualizações financeiras, Helena conseguirá pagar suas despesas com o novo salário? Caso falte ou sobre dinheiro, indique o valor exato da diferença em reais.
\par\vspace{\espacoQuestao}

\noindent\textcolor{CorLinha}{\rule{\linewidth}{0.5pt}} 
\par\vspace{\espacoPreQuestao}
\textbf{44. (Estilo UNESP - Refutação de Distratores)} - Programas de edição gráfica alteram as dimensões de uma imagem usando proporções percentuais, o que afeta severamente a área total de impressão. Um designer possui uma fotografia retangular cujas dimensões originais são \textbf{20 cm} de comprimento por \textbf{15 cm} de largura. 

Para adequar a imagem a um site, ele aplica uma redução geométrica de \textbf{20\%} no comprimento e, simultaneamente, \textbf{20\%} na largura. Qual foi a perda percentual real da \textbf{área} total da fotografia?
\begin{enumerate}[label=\alph*)]
    \item 20\%
    \item 36\%
    \item 40\%
    \item 64\%
    \item 100\%
\end{enumerate}
\begin{boxresponda}
\textbf{Responda:} Assinale a alternativa correta demonstrando o cálculo das áreas. Em seguida, justifique por que as alternativas \textbf{a)} e \textbf{c)} são armadilhas lógicas comuns em provas de concursos.
\end{boxresponda}
\par\vspace{\espacoQuestao}

\noindent\textcolor{CorLinha}{\rule{\linewidth}{0.5pt}} 
\par\vspace{\espacoPreQuestao}
\textbf{45.} Durante uma eleição para o grêmio estudantil de um colégio, três chapas disputaram os votos dos alunos. O regulamento previa que os votos nulos e brancos não seriam contabilizados para a definição dos percentuais \textbf{válidos}. 

Apurados os \textbf{800} votos totais depositados nas urnas, verificou-se que \textbf{15\%} correspondiam a votos brancos ou nulos. Sabendo que a Chapa A recebeu \textbf{40\%} dos votos válidos, qual foi o número absoluto de alunos que votou na Chapa A?
\par\vspace{\espacoQuestao}

\noindent\textcolor{CorLinha}{\rule{\linewidth}{0.5pt}} 
\par\vspace{\espacoPreQuestao}
\textbf{46.} A diferença conceitual entre margem de lucro sobre o custo e margem de lucro sobre a venda é uma métrica vital da contabilidade. Uma fábrica produz camisetas esportivas a um custo de fabricação de \textbf{R\$ 40,00} por unidade. Essas mesmas camisetas são repassadas para as lojas e vendidas ao consumidor final por \textbf{R\$ 50,00}.

Calcule e responda de forma estruturada: qual é a taxa percentual de lucro obtida utilizando como base o preço de \textbf{custo}? E qual é a taxa percentual de lucro caso o setor financeiro decida utilizar como base o preço de \textbf{venda}?
\par\vspace{\espacoQuestao}

\noindent\textcolor{CorLinha}{\rule{\linewidth}{0.5pt}} 
\par\vspace{\espacoPreQuestao}
\textbf{47.} Na pesquisa de satisfação de um refeitório corporativo, diagramas lógicos são utilizados para mapear a aceitação do cardápio. O diagrama de Venn exibe o perfil percentual de preferência entre as opções de refeição.

\begin{center}
\begin{adjustbox}{max width=\linewidth}
\begin{tikzpicture}
\definecolor{fundo}{HTML}{F4F6F7}
\definecolor{circA}{HTML}{E74C3C} % Carne
\definecolor{circB}{HTML}{3498DB} % Frango
\definecolor{texto}{HTML}{2C3E50}

% Sombra Projetada
\fill[black!15, rounded corners=4pt, drop shadow={opacity=0.2, shadow xshift=3pt, shadow yshift=-3pt}] (-3, -3) rectangle (7, 4);
\fill[fundo, rounded corners=4pt] (-3, -3) rectangle (7, 4);
\draw[thick, black!70, rounded corners=4pt] (-3, -3) rectangle (7, 4);

% Universo
\node[anchor=north west, font=\bfseries] at (-2.8, 3.8) {Total Entrevistado: 100\%};
\node[anchor=south east, font=\bfseries, fill=fundo, inner sep=2pt] at (6.8, -2.8) {Nenhum dos dois: 15\%};

% Círculos com opacidade
\fill[circA, opacity=0.3] (0, 0) circle (2.2);
\draw[thick, circA] (0, 0) circle (2.2);

\fill[circB, opacity=0.3] (3, 0) circle (2.2);
\draw[thick, circB] (3, 0) circle (2.2);

% Rótulos com fill=white
\node[font=\bfseries, text=texto, fill=white, inner sep=2pt, rounded corners=1pt] at (0, 2.5) {Carne (60\%)};
\node[font=\bfseries, text=texto, fill=white, inner sep=2pt, rounded corners=1pt] at (3, 2.5) {Frango (55\%)};

% Valor na intersecção
\node[font=\bfseries, text=black, fill=white, inner sep=2pt, rounded corners=1pt] at (1.5, 0) {x\%};
\end{tikzpicture}
\end{adjustbox}
\end{center}

Baseando-se nos dados apresentados, equacione o problema da intersecção de conjuntos e determine a taxa percentual de funcionários (\textbf{x\%}) que declararam gostar simultaneamente das duas opções de proteína (Carne e Frango).
\par\vspace{\espacoQuestao}

\noindent\textcolor{CorLinha}{\rule{\linewidth}{0.5pt}} 
\par\vspace{\espacoPreQuestao}
\textbf{48.} Operações volumétricas costumam esconder armadilhas com casas decimais. Uma caixa d'água em formato de paralelepípedo retangular reto possui dimensões internas de \textbf{1,5 m} de comprimento, \textbf{1,2 m} de largura e \textbf{0,8 m} de altura.

Sabendo que a fórmula geométrica do volume é o produto das três dimensões (comprimento $\times$ largura $\times$ altura), calcule o volume exato da caixa d'água em metros cúbicos ($ \text{m}^3 $).
\par\vspace{\espacoQuestao}

\noindent\textcolor{CorLinha}{\rule{\linewidth}{0.5pt}} 
\par\vspace{\espacoPreQuestao}
\textbf{49.} Uma jarra contém \textbf{2,5 L} de suco de laranja concentrado. Para tornar o suco mais suave e adequado para uma festa, a dona de casa decide acrescentar exatamente \textbf{35\%} de água mineral ao volume já existente na jarra.

Determine, em forma de número decimal, o volume absoluto de água adicionado e a nova capacidade total (em litros) de líquido contido na jarra.
\par\vspace{\espacoQuestao}

\noindent\textcolor{CorLinha}{\rule{\linewidth}{0.5pt}} 
\par\vspace{\espacoPreQuestao}
\textbf{50.} No laboratório de geologia, uma amostra mineral em formato irregular foi colocada em uma balança de precisão, registrando \textbf{14,75 g}. Durante os testes químicos de pureza, essa rocha foi submetida a um forte ácido e perdeu \textbf{20\%} de sua massa original devido à corrosão.

Apresente a operação que comprova a perda mássica e registre o peso final da amostra de rocha.
\par\vspace{\espacoQuestao}

\noindent\textcolor{CorLinha}{\rule{\linewidth}{0.5pt}} 
\par\vspace{\espacoPreQuestao}
\textbf{51. (Estilo UERJ)} - A conservação de alimentos por desidratação (secagem) altera drasticamente a concentração de massa dos compostos. Um fazendeiro colheu \textbf{100 kg} de melancias maduras. Análises laboratoriais mostram que \textbf{99\%} da massa física dessa fruta fresca é composta estritamente por água.

Após algumas semanas armazenadas, as melancias perderam parte dessa umidade por evaporação, de modo que a água passou a representar exatamente \textbf{98\%} da massa total das frutas. 
\begin{boxresponda}
\textbf{Responda:} Sabendo que apenas a água evaporou e que a "massa sólida" da melancia permaneceu intacta, calcule a nova massa total (em kg) desse lote. Explique detalhadamente o fenômeno de conservação de bases que resolve este paradoxo matemático.
\end{boxresponda}
\par\vspace{\espacoQuestao}

\noindent\textcolor{CorLinha}{\rule{\linewidth}{0.5pt}} 
\par\vspace{\espacoPreQuestao}
\textbf{52. (ENEM)} - Um produtor de café empacota seus produtos em fardos idênticos pesando \textbf{15,5 kg} cada. Uma distribuidora encomendou um pedido equivalente a \textbf{465 kg} desse café. Por questões de fidelidade comercial, o produtor concedeu um desconto de \textbf{R\$ 0,50} por cada quilograma faturado.

Sabendo que o preço normal do quilograma do café é \textbf{R\$ 18,50}, determine a quantidade de fardos que o caminhão deverá transportar e o valor financeiro final (já com os descontos) pago pela distribuidora.
\par\vspace{\espacoQuestao}

\end{multicols*}
\end{document}